\section{Related work}
\label{sec:related}



\subsection{Graph statistics and model performance}

Edge homophily summarizes whether adjacent nodes share labels, but a global fraction omits class imbalance, feature geometry, neighborhood distributions, and local variation~\citep{mcpherson2001homophily,zhu2020beyond,ma2021homophily}. Post-aggregation similarity, adjusted homophily, label informativeness, and local characterizations describe additional aspects of message-passing behavior~\citep{luan2022linkx,platonov2023characterizing,loveland2024local}.

Several diagnostics address relative model performance directly. \citet{luan2023whendo} analyze node distinguishability and introduce the Classifier-based Performance Metric (CPM). \citet{zheng2024trihom} disentangle label, structural, and feature homophily; their 31-dataset study compares Tri-Hom with 17 metrics, including correlations with GCN-, GraphSAGE-, and GAT-versus-MLP gaps. Under a heterophilous stochastic block model, \citet{pmlr-v235-wang24u} relate separability gain to neighborhood-distribution distance and degree. These approaches characterize richer signals than the simple threshold rules used in our primary case studies. The post-hoc CPM-GNB adaptation in Section~\ref{sec:cpm_results} tests one such feature-based signal under the same decision audit.

Benchmark protocols affect the interpretation of these comparisons. Repeated splits and standardized graph collections expose sensitivity to splitting, tuning, and dataset selection~\citep{shchur2018pitfalls,hu2020ogb}. Reassessments of heterophilous benchmarks and homophily metrics further emphasize protocol quality~\citep{platonov2023heterophilous,luan2024reevaluating}. Dataset-level inference distinguishes variation across tasks from repeated seeds on the same task~\citep{demsar2006statistical}. We use paired evaluation and restrict label-dependent diagnostics to training labels. Recent preprint work also examines homophily-aware split construction to control relational composition across folds~\citep{jami2026stratification}. Our comparisons use the retained splits; sensitivity to alternative split designs remains untested.

\subsection{Architectures and the candidate portfolio}

Graph architectures differ in the signals they propagate and combine. H2GCN separates ego and neighbor representations and uses two-hop neighborhoods~\citep{zhu2020beyond}; LINKX separates structural and feature processing~\citep{lim2021linkx}; GPR-GNN learns a combination of propagation depths~\citep{chien2021gprgnn}. A recent preprint uses final-layer and head refitting to separate recoverable information from classification failure~\citep{ness2026rarity}. This distinction reinforces the need to audit the trained baseline before attributing a model gap to graph structure. Training depth, regularization, and representation contraction also affect predictive performance~\citep{oono2020exponential,chen2020measuring,rong2020dropedge,zhao2020pairnorm}. A one-hop mixing statistic therefore faces a different prediction problem when the graph action selects among these trained candidates. We examine that dependence using GCN, GAT, GraphSAGE, H2GCN, LINKX, and GPR-GNN, then restrict the portfolio while holding the diagnostic decisions fixed.

\subsection{Model selection and graph-signal controls}

Graph model selection already includes architecture search and learned selection across tasks. Auto-GNN searches architectures for a graph task~\citep{zhou2020autognn}. MetaGL selects graph-learning models using meta-graph features and historical performance~\citep{park2023metagl}; GLEMOS benchmarks instantaneous selection across 366 models and 457 graphs~\citep{park2023glemos}.

Using a published statistic for model selection requires specifying its label access, graph operator, input map, and score-to-action rule. A correlation study alone does not specify these choices.

These approaches predict different outcomes. CPM uses feature-based hypothesis tests to indicate the advantage of a graph-aware model over its paired graph-agnostic model~\citep{luan2023whendo}. Tri-Hom combines label, structural, and feature information; its evaluation measures correlations with model performance and graph--MLP gaps~\citep{zheng2024trihom}. GLEMOS selects models on unseen graphs from historical performance records and graph features, before fitting candidates on those graphs. It evaluates selection and ranking quality across transfer settings~\citep{park2023glemos}. We examine full-policy accuracy loss, abstention fallback, and available improvement as the candidate portfolio changes. Fixed threshold policies allow each decision to be traced, although they are not representative of all published diagnostics. The analysis uses established evaluation principles to locate policy losses and reversals in relative performance.

Our target is a binary graph-portfolio-versus-MLP action evaluated through oracle-portfolio regret and abstention coverage. The reference policies are always-graph, always-MLP, random, and direct validation selection. Each architecture receives four trials, yielding unequal aggregate search budgets for the six-architecture graph action and the single-architecture MLP action. This resource asymmetry is part of the evaluated decision problem.

Strong feature-only baselines and edge controls help interpret that decision. Well-tuned MLPs can be competitive when node features encode graph information~\citep{katsman2026necessary}, while edge-removal controls show that models can exploit graph structure even when it hurts generalization~\citep{bechler2024graphs}. Our degree-preserving rewiring intervention provides complementary evidence about graph signal. It changes class mixing and higher-order organization together, so it does not isolate homophily. Its role is to assess structural signal separately from the utility of the diagnostic policies.

The fixed-degree Gaussian calculation in Appendix~\ref{sec:snr} describes how correlation, degree, and feature strength interact under a prescribed linear average. Contextual stochastic-block-model recovery theory and finite-depth embedding analyses provide related theoretical perspectives~\citep{lu2023contextual,pmlr-v269-dalle25a}. The calculation describes aggregation moments; the recorded policy comparisons measure decision utility within the evaluated configurations.
